\subsection{Binary Access Patterns}

Binary access usually cares less about lines and more about byte positions, byte counts, and raw transfer. The program may read a fixed number of bytes, write prepared byte buffers, or move to a specific byte offset inside the file.

This is different from text access. Text mode thinks in decoded characters and newline translation. Binary mode thinks in byte sequences.

\subsubsection{Reading Fixed-Size Chunks for Large Files}

The \texttt{read()} method accepts an optional size argument. In binary mode, that size means number of bytes. This allows a program to process a file gradually instead of loading the whole file into memory.

\begin{verbatim}
path = "chunk_read_demo.bin"

data = (
    b"Chunk 1: D'oh!\n"
    b"Chunk 2: Serenity now!\n"
    b"Chunk 3: The answer is 42.\n"
    b"Chunk 4: Nobody expects the Spanish Inquisition!\n"
)

with open(path, "wb") as file_obj:
    file_obj.write(data)

chunk_size = 12

with open(path, "rb") as file_obj:
    print("Binary file object:")
    print(f"file_obj       = {file_obj}")
    print(f"type(file_obj) = {type(file_obj)}")

    print()
    print("Selected names from dir(file_obj):")
    for name in ["read", "write", "seek", "tell", "close", "mode"]:
        print(f"{name:10} -> {name in dir(file_obj)}")

    print()
    print("Reading fixed-size chunks:")

    chunk_number = 1

    while True:
        chunk = file_obj.read(chunk_size)

        if chunk == b"":
            break

        print(f"chunk {chunk_number:02d}: {chunk!r}")
        print(f"    type(chunk) = {type(chunk)}")
        print(f"    len(chunk)  = {len(chunk)}")

        chunk_number += 1
\end{verbatim}

The loop ends when \texttt{read(chunk\_size)} returns the empty byte string \texttt{b""}. That is the binary-mode end-of-file marker for this pattern.

The important memory rule is that only one chunk is held at a time by the loop variable \texttt{chunk}. This is the normal pattern for large binary files, because the program does not need to materialize the whole file as one large \texttt{bytes} object.

\subsubsection{Writing Byte Buffers to External Storage}

Binary writing sends byte-oriented objects to a file opened in binary write mode. The most common objects are \texttt{bytes} and \texttt{bytearray}.

\begin{verbatim}
path = "byte_buffer_write_demo.bin"

header = b"HEADER:"
body = bytearray(b" fact number 41")
body[-1] = ord("2")
ending = b"\n"

print("Objects before writing:")
print(f"header       = {header}")
print(f"type(header) = {type(header)}")

print(f"body         = {body}")
print(f"type(body)   = {type(body)}")

print(f"ending       = {ending}")
print(f"type(ending) = {type(ending)}")

print()
print("Selected names from dir(body):")
for name in ["append", "extend", "insert", "pop", "decode", "hex"]:
    print(f"{name:10} -> {name in dir(body)}")

with open(path, "wb") as file_obj:
    count1 = file_obj.write(header)
    count2 = file_obj.write(body)
    count3 = file_obj.write(ending)

print()
print("Byte counts returned by write():")
print(f"count1 = {count1}")
print(f"count2 = {count2}")
print(f"count3 = {count3}")

with open(path, "rb") as file_obj:
    restored_data = file_obj.read()

print()
print("Restored byte data:")
print(f"restored_data       = {restored_data}")
print(f"type(restored_data) = {type(restored_data)}")

print()
print("Decoded only for display:")
print(restored_data.decode("utf-8"))
\end{verbatim}

The file receives bytes. It does not receive Python strings, and it does not perform text encoding by itself. The \texttt{bytearray} object is accepted because it is a bytes-like object.

A list of byte chunks can also be written one by one.

\begin{verbatim}
path = "many_buffers_demo.bin"

chunks = [
    b"First block: Spam.\n",
    b"Second block: Eggs.\n",
    b"Third block: 42.\n",
]

print("Chunk container:")
print(f"chunks       = {chunks}")
print(f"type(chunks) = {type(chunks)}")
print(f"len(chunks)  = {len(chunks)}")

print()
print("Selected names from dir(chunks):")
for name in ["append", "extend", "insert", "pop", "sort", "reverse"]:
    print(f"{name:10} -> {name in dir(chunks)}")

with open(path, "wb") as file_obj:
    total_count = 0

    for chunk in chunks:
        written_count = file_obj.write(chunk)
        total_count += written_count

        print(f"wrote {written_count:02d} bytes from {chunk!r}")

print()
print(f"total_count = {total_count}")

with open(path, "rb") as file_obj:
    data = file_obj.read()

print()
print(data)
\end{verbatim}

This pattern is useful when a binary output is assembled from several pieces: a header, a payload, metadata, padding bytes, or generated byte buffers.

\subsubsection{Random Access with \texttt{.seek()} and \texttt{.tell()}}

Sequential reading consumes bytes in order. Random access means moving to a specific byte position and reading or writing from there.

The \texttt{tell()} method reports the current byte position. The \texttt{seek()} method moves the current position.

\begin{verbatim}
path = "seek_tell_demo.bin"

data = b"0123456789ABCDEFGHIJKLMNOPQRSTUVWXYZ"

with open(path, "wb") as file_obj:
    file_obj.write(data)

with open(path, "rb") as file_obj:
    print("Initial position:")
    print(f"file_obj.tell() = {file_obj.tell()}")

    part1 = file_obj.read(5)

    print()
    print("After reading 5 bytes:")
    print(f"part1           = {part1}")
    print(f"type(part1)     = {type(part1)}")
    print(f"file_obj.tell() = {file_obj.tell()}")

    file_obj.seek(10)

    print()
    print("After seek(10):")
    print(f"file_obj.tell() = {file_obj.tell()}")

    part2 = file_obj.read(6)

    print()
    print("Read 6 bytes from position 10:")
    print(f"part2           = {part2}")
    print(f"file_obj.tell() = {file_obj.tell()}")

    file_obj.seek(0)

    print()
    print("After seek(0), back to start:")
    print(f"file_obj.tell() = {file_obj.tell()}")

    part3 = file_obj.read(3)

    print()
    print("Read first 3 bytes again:")
    print(f"part3           = {part3}")
    print(f"file_obj.tell() = {file_obj.tell()}")
\end{verbatim}

The byte positions are zero-based. Position \texttt{0} is the first byte. After reading five bytes, the current position is \texttt{5}, meaning the next read begins at byte index \texttt{5}.

Random access also works for overwriting bytes in an existing binary file.

\begin{verbatim}
path = "binary_overwrite_demo.bin"

with open(path, "wb") as file_obj:
    file_obj.write(b"Answer: 41\n")

with open(path, "r+b") as file_obj:
    print("Before overwrite:")
    file_obj.seek(0)
    before_data = file_obj.read()
    print(before_data)

    file_obj.seek(9)
    file_obj.write(b"2")

    print()
    print("After overwrite:")
    file_obj.seek(0)
    after_data = file_obj.read()
    print(after_data)

print()
print("Decoded for display:")
print(after_data.decode("utf-8"))
\end{verbatim}

The mode \texttt{"r+b"} opens the file for both reading and writing in binary mode without erasing the existing file. The call \texttt{seek(9)} moves to the byte position containing \texttt{b"1"}. The write then replaces that one byte with \texttt{b"2"}.

This is not line-based editing. It is byte-position editing. If the replacement has a different length, the surrounding bytes are not automatically shifted like elements in a Python list. Binary random access should therefore be used with formats where byte positions and byte sizes are understood.

The practical binary-access rules are:

\begin{itemize}
    \item use \texttt{read(size)} to process large files in bounded chunks;
    \item use \texttt{write()} with \texttt{bytes} or \texttt{bytearray} buffers;
    \item use \texttt{tell()} to inspect the current byte position;
    \item use \texttt{seek()} to move to a specific byte position.
\end{itemize}

